% Fig. 2: node connections as a single-line block diagram (single column).
% Source: plan.md (Hardware and Measurement Design) and research.md R-001/R-002.
% Pin numbers are omitted: they are candidates until the PCB review.
% Replace or complement with a labelled photo once the prototype is built.
\begin{tikzpicture}[
  font=\scriptsize,
  every node/.append style={execute at begin node={\hyphenpenalty=10000\exhyphenpenalty=10000}},
  blk/.style={draw, rounded corners=1pt, fill=gray!12, align=center,
    minimum height=7mm, inner sep=2pt},
  mains/.style={blk, fill=white, text width=18mm, minimum height=6mm},
  iso/.style={blk, fill=white, text width=16mm, double, double distance=0.6pt},
  fuse/.style={draw, fill=white, minimum width=4mm, minimum height=2mm, inner sep=0pt},
  wire/.style={line width=1pt},
  arr/.style={-{Latex[length=1.5mm]}, semithick},
  bus/.style={{Latex[length=1.5mm]}-{Latex[length=1.5mm]}, semithick},
  lbl/.style={font=\tiny, inner sep=1pt, align=center},
  ttl/.style={font=\scriptsize\bfseries, align=center},
]
% --- Mains side (single-line) ---
\node[mains] (grid) at (0.9,0.95) {Red monofásica\\220 V, 60 Hz};
\node[mains] (rcd) at (0.9,0.1) {Diferencial 30 mA};
\node[mains] (ocp) at (0.9,-0.75) {Protección bipolar 10 A};
\node[mains] (load) at (0.9,-4.75) {Carga de prueba};
\draw[wire] (grid) -- (rcd);
\draw[wire] (rcd) -- (ocp);
\draw[wire] (ocp) -- (load);

% --- Isolating devices (straddle the barrier) ---
\node[iso] (psu) at (3.35,-1.6) {Fuente aislada certificada};
\node[iso] (vs) at (3.35,-2.7) {Sensado de tensión aislado};
\node[iso] (ct) at (3.35,-3.8) {SCT-013-020 (20 A/1 V)};
\draw[semithick] (0.9,-1.6) -- node[fuse, pos=0.55] {} (psu.west);
\draw[semithick] (0.9,-2.7) -- node[fuse, pos=0.55] {} (vs.west);
\node[lbl, above] at (1.85,-2.6) {fusible};
\draw[semithick] (0.9,-3.8) ellipse (2.2mm and 1.1mm);
\draw[semithick, densely dashed] (1.12,-3.8) -- node[lbl, above] {un conductor} (ct.west);
\fill (0.9,-1.6) circle (0.8pt) (0.9,-2.7) circle (0.8pt);

% --- Low-voltage side ---
\node[blk, minimum width=17mm] (mcu) at (5.55,-1.6) {ESP32};
\node[blk, minimum width=17mm] (adc) at (5.55,-2.7) {ADS131M02\\2 canales};
\node[blk, text width=10mm] (led) at (7.75,-1.6) {LED o\\zumbador};
\node[blk, text width=10mm] (mem) at (7.75,-2.7) {FRAM\\y RTC};
\draw[arr] (psu) -- node[lbl, above] {5 V} (mcu);
\draw[arr] (vs) -- node[lbl, below] {canal V} (adc);
\draw[arr] (ct.east) -| node[lbl, pos=0.25, above] {canal I} (adc.south);
\draw[bus] (adc) -- node[lbl, right] {SPI, DRDY} (mcu);
\draw[arr] (mcu.east |- led.west) -- node[lbl, above] {GPIO} (led.west);
\draw[bus] ($(mcu.east)+(0,-0.22)$) -| (6.75,-2.7) -- (mem.west);
\node[lbl, right] at (6.75,-2.2) {I\textsuperscript{2}C};
\draw[arr] (mcu.north) -- ++(0,0.45) node[lbl, above] {Wi-Fi hacia la red del laboratorio};

% --- Barrier and titles ---
\begin{scope}[on background layer]
  \draw[very thick, dash pattern=on 3pt off 2pt] (3.35,1.25) -- (3.35,-4.55);
\end{scope}
\node[lbl, fill=white] at (3.35,-4.75) {barrera de aislamiento};
\node[ttl, anchor=south] at (0.9,1.4) {Lado de red};
\node[ttl, anchor=south] at (6.0,1.4) {Baja tensión (envolvente)};
\end{tikzpicture}
